\section{Introduction}\label{a:latex_en}

\subsection{Motivation}

Particle accelerators are one of the most advanced scientific devices ever developed 
to perform innovative scientific experiments. The Deutsches Elektronen-Synchrotron 
(DESY) \cite{DESYwebsite} located in Hamburg region hosts multiple top notch accelerator facilities such as 
European XFEL X-ray Laser and the PETRA III. These facilities accelerate charged particles, that enable researchers to study atomic level processes and structures with 
great precision. A crucial challenge for particle accelerators is the tuning process, which 
involves adjustment of many magnetic components to produce desired beam properties including 
position, size and divergence. As a specialized test facility for accelerators research and 
development, the ARES (Accelerator Research Experiment at SINBAD), which is a linear accelerator at DESY,
allows the evaluation of innovative tuning techniques prior to deployment on even larger facilities. \

Beam tuning is essentially an optimization task. In order to find exact configurations that minimize the 
deviation between the actual beam properties and the desired beam, we optimize the adjustable parameters 
(like magnet strengths and corrector/steering angles). However, the optimization is not that simple and 
perhaps contains a number of challenges: 

\begin{enumerate}
    \item \textbf{Expensive Evaluation:} It is impractical to conduct an extensive search since measurement of each beam property takes time and resources.
    \item \textbf{Uncertainity:} There are a number of systematic errors and noises that might affect beam measurement.
    \item \textbf{Black-box Nature:} The complex relationship between input parameters and beam characteristics cannot be directly observed analytically during operation.
    \item \textbf{Unknown / Partial Observability:} Important parameters like characteristics of incoming beam and magnet misalignments are frequently either unknown or partially observed. 

\end{enumerate}

A traditional approach to accelerator tuning could be through skilled professionals who manually adjust 
parameters using their knowledge and experience. While being efficient, this method significantly relies 
on individual skill, takes a lot of time, is expensive and does not scale well to more complicated accelerator 
systems. Therefore, automated optimization methods are necessary for effective accelerator performance.
Several optimization algorithms have been applied to particle accelerator tuning problems with their own unique 
advantages and disadvantages. Nelder-Mead Simplex is a gradient free optimization technique and one of the approaches 
that is widely used because of its simplicity and robustness \cite{nelder1965simplex}. However they do not create models that generalize across many 
optimization runs and usually need many function evaluations to converge. Another approach is the Reinforcement Learning (RL) 
technique that has shown remarkable performance \cite{kaiser2024rl}. RL agents can learn to tune accelerators in very few steps during deployment 
by training neural network regulations in simulation. However, RL requires extensive training data as well as very careful reward 
function design. The trained policies could not generalize well to conditions outside of their training distribution. The approach 
of our interest for this project work is Bayesian Optimization (BO), which finds a middle ground among these approaches. BO creates 
a probabilistic surrogate model of the objective function, which is usually a Gaussian Process, that provides estimates of uncertainty 
as well as predictions. This makes it possible for the assessment points to wisely balance between both exploration and exploitation. BO is 
a desirable option for accelerator tuning since it excels at solving expensive black-box optimization problems. \

Another key insight motivating this project work is that, since the physics of particle accelerators is already well understood now, 
We can map the passage of charged particles via magnets using a reliable model and then they travel through magnetic fields according to 
well established physics knowledge. Thus, all thanks to modern simulation library like Cheetah \cite{kaiser2024bridging} that we can now predict beam 
behavior with high accuracy. The standard BO uses uninformed zero mean prior function in the Gaussian process model. This implies that 
the algorithm must learn the whole function landscape from scratch and have to begin with no prior knowledge of the intended objective 
function values. On the other hand, if the physics informed knowledge could be incorporated into the prior, this shall give advantage 
and optimization could already start with meaningful predictions that could focus its limited evaluation budget on identifying the 
differences between the model and the reality. The importance of physics informed prior for accelerator optimization has also been 
shown in the recent work \cite{boltz2025priormeanmodule} where the neural network surrogate models trained on the simulation data can perform as effective prior 
mean functions, significantly speeding up the convergence. However, training neural network surrogates requires very large datasets 
from simulations. So let's continue studying Bayesian Optimization with Cheetah informed prior.




\subsection{Extension from Simple FODO to Complex ARES Problem}

Cheetah has already been demonstrated as a physics informed prior for BO on a very simple and straightforward beam focusing FODO 
(Focus-Drift-Defocus-Drift) lattice \cite{boltz2025priormeanmodule}. A FODO cell is a fundamental building block made up of alternating quadrupole magnets 
for focusing and defocusing that are separated by drift spaces. \

Since this optimization problem was two dimensional, thus it only involved strengths of two quadrupole magnets ($k^{Q1}$ and $k^{Q2}$) with 
the objective of minimising beam size at target location, defined in equation \ref{mae_fodo}. Here $\sigma_x$ and $\sigma_y$ are the horizontal 
and vertical beam sizes, considering only beam size (focusing) without beam position (centering), unlike ARES problem, 
refelecting the simpler nature of FODO demonstration.
\begin{equation}
\label{mae_fodo}
mae = 0.5(|\sigma_x|+|\sigma_y|)
\end{equation}

Moreover, although there were two drift spaces ($d_1$ and $d_2$), they were constrained to have same lengths between the quadrupoles, 
therefore only a single trainable hyperparameter was there, which was restricted to the interval [0.2, 1.0] m and was included in the physics informed 
prior mean function. The results showed that BO with Cheetah prior converged about 15 steps, indicating better sample efficiency 
than Nelder Mead Simplex and standard BO with zero prior. Furthermore, the method worked even with a mismatched prior, the algorithm 
learned the correct values during GP model fitting and converged to the true optimum when the drift lengths in the prior model (0.5m)
was different from the ground truth (0.7m). 
While this validated the fundamental concept, the FODO problem is an oversimplified example which does not depict the complex nature 
of actual accelerator tuning operations. The two dimensional search space is quite easy to explore and the model reality mismatch was 
restricted to a single geometric parameter (drift length) instead of the various other sources of uncertainty found in actual accelerators. \

This project extends the FODO proof of concept to a significantly more complex and realistic problem which is transverse beam tuning at the 
ARES Experimental Area at DESY. The extension contain following key advancements: \

\begin{enumerate}
    \item \textbf{Higher Dimensionality:} In contrast to two dimensional FODO problem, ARES problem possess five dimensional search space including two steering/corrector magnet angles ($\theta^{CV}$ and $\theta^{CH}$) and three quadrupole strengths ($k^{Q1},k^{Q2}, k^{Q3}$). The volume of the search space is greatly increased by around 2.5x in dimensionality, which makes it more difficult and raises the importance of informed prior. 
    \item \textbf{Realistic Lattice:} The ARES Experimental Area is a real accelerator section loaded from an actual lattice description file unlike the synthetic FODO cell. 
    \item \textbf{Objective Nature:} For Ares problem, the objective function combines four beam parameters into a single metric, including horizontal and vertical sizes ($\sigma_x$, $\sigma_y$) as well as the positions ($\mu_x$, $\mu_y$). This makes simultaneous optimization of beam centering and focusing.
    \item \textbf{Trainable Quadrupole's Misalignments:} Under this project work, the use of quadrupole misalignments as the cause of model reality mismatch is an important step forward. The ARES prior includes 6 trainable misalignment hyperparameters (horizontal and vertical for each quadrupoles Q1, Q2 and Q3) which are limited to physically plausible intervals of $\pm 0.5 mm$, in contrast to the FODO demonstration, where only a single hyperparameter (drift length) was learned.
\end{enumerate}


\subsection{Contribution}

Thus the main contributions of this work are : \

\begin{enumerate}
    \item Demonstrating that Cheetah informed prior approach is scalable from 2D synthetic lattice to 5D real world accelerator tuning problem. 
    \item Implementing physics informed prior mean function for the ARES lattice being compatible with the Xopt Bayesian optimisation framework
    \item Extending prior with six trainable misalignment parameters, enabling the physics model to adapt to real accelerator conditions during optimization
    \item Containing a comprehensive comparison between standard BO, Nealder-Mead and physic-informed BO (matched and mismatched prior) across multiple trials.
\end{enumerate}